
\documentclass[10pt]{article}
\usepackage[utf8]{inputenc}
\usepackage[T1]{fontenc}
\usepackage{amsmath, amssymb, xcolor, geometry, enumitem, multicol, tcolorbox}

\definecolor{mainblue}{RGB}{0, 102, 204}
\geometry{a4paper, margin=1.2cm, top=1cm, bottom=3cm, footskip=1.2cm}

\begin{document}
\enlargethispage{2.5cm}

% Modern Header
\begin{tcolorbox}[colback=mainblue!5, colframe=mainblue, sharp corners, boxrule=0.5pt]
    \small \textbf{Unit:} undefined \hfill \textbf{Name:} \rule{3cm}{0.4pt} \\
    \small \textbf{Topic:} Variables and Expressions / Order of Operations \hfill \textbf{Date:} \rule{2cm}{0.4pt} \\
    \centering \textbf{\large Foundation (Level -1)}
\end{tcolorbox}

\vspace{0.2cm}

% MCQ Section
\begin{multicols}{2}
\begin{enumerate}[label=\textbf{Q\arabic*.}, leftmargin=*, itemsep=6pt, parsep=0pt]

  \item \small What is the variable in the expression $2x$?
  \begin{enumerate}[label=(\Alph*), leftmargin=0.6cm, nosep, itemsep=2pt]
    \item $2$
    \item $x$
    \item $2x$
    \item None of the above
    \item Both A and B
  \end{enumerate}

  \item \small Identify the coefficient in the term $4y$.
  \begin{enumerate}[label=(\Alph*), leftmargin=0.6cm, nosep, itemsep=2pt]
    \item $y$
    \item $4$
    \item $4y$
    \item None of the above
    \item Both A and B
  \end{enumerate}

  \item \small What is the operator in the expression $3 + 2$?
  \begin{enumerate}[label=(\Alph*), leftmargin=0.6cm, nosep, itemsep=2pt]
    \item $3$
    \item $2$
    \item $+$
    \item None of the above
    \item Both A and B
  \end{enumerate}

  \item \small In the expression $x - 1$, what is being subtracted?
  \begin{enumerate}[label=(\Alph*), leftmargin=0.6cm, nosep, itemsep=2pt]
    \item $x$
    \item $1$
    \item $-$
    \item None of the above
    \item Both A and B
  \end{enumerate}

  \item \small What is the constant term in the expression $2x + 3$?
  \begin{enumerate}[label=(\Alph*), leftmargin=0.6cm, nosep, itemsep=2pt]
    \item $2x$
    \item $3$
    \item $2x + 3$
    \item None of the above
    \item Both A and B
  \end{enumerate}

  \item \small Identify the variable in $5z$.
  \begin{enumerate}[label=(\Alph*), leftmargin=0.6cm, nosep, itemsep=2pt]
    \item $5$
    \item $z$
    \item $5z$
    \item None of the above
    \item Both A and B
  \end{enumerate}

  \item \small What operation is represented by the symbol $*$?
  \begin{enumerate}[label=(\Alph*), leftmargin=0.6cm, nosep, itemsep=2pt]
    \item Addition
    \item Subtraction
    \item Multiplication
    \item Division
    \item None of the above
  \end{enumerate}

  \item \small In the expression $7 - x$, what is being subtracted from $7$?
  \begin{enumerate}[label=(\Alph*), leftmargin=0.6cm, nosep, itemsep=2pt]
    \item $7$
    \item $x$
    \item $-$
    \item None of the above
    \item Both A and B
  \end{enumerate}

\end{enumerate}
\end{multicols}

\vspace{-0.2cm}
\hrule
\vspace{0.2cm}
\textbf{\small Open Ended Questions (FRQ)}
\begin{enumerate}[label=\textbf{Q\arabic*.}, start=9, itemsep=5pt, topsep=0pt]

  \item \small Draw a simple diagram to represent the expression $2x + 1$. Explain your drawing in one or two sentences. \vspace{1.6cm}

  \item \small Explain in your own words what a variable is in an algebraic expression. Provide an example. \vspace{1.6cm}

\end{enumerate}

\vfill

% Chef's Tips Footer
\begin{tcolorbox}[colback=blue!5, colframe=mainblue!50, title=\small \textbf{Chef's Tips}, fonttitle=\bfseries, bottom=1mm]
\begin{itemize}[noitemsep, topsep=2pt, leftmargin=0.5cm]
  \item \scriptsize Focus on basic conceptual definitions and single-step identification to ensure a strong foundation in variables and expressions.\item \scriptsize Use visual aids like number lines, graphs, or simple diagrams to help explain abstract concepts like variables and constants.\item \scriptsize Encourage students to use their own words to describe algebraic concepts, promoting a deeper understanding of the material.
\end{itemize}
\end{tcolorbox}

\end{document}
  